% APROVEITAMENTO E VALIDAÇÃO DO CONHECIMENTO ----------------------------------------------

\clearpage
\section{APROVEITAMENTO E VALIDAÇÃO DO CONHECIMENTO}

No que se refere ao aproveitamento de componentes curriculares cursados, o IFCE assegurará aos estudantes ingressantes e veteranos o direito de aproveitamento, mediante análise, desde que haja compatibilidade de programa e carga horária de, no mínimo, 75\% do total estipulado para o componente curricular a ser aproveitado. O componente curricular apresentado deve estar no mesmo nível de ensino ou em um nível de ensino superior ao do componente curricular a ser aproveitado, devendo ser solicitado no máximo uma vez. 

No aproveitamento, deverão ser considerados os conhecimentos adquiridos não só para as disciplinas do semestre em curso, como também para as de semestres posteriores, no caso de aluno recém-ingresso. Este, terá 10 (dez) dias após a sua matrícula, para requerer o aproveitamento de disciplina. Quanto ao aluno veterano, o aproveitamento será para o semestre/ano posterior, devendo a solicitação ser feita durante os 30 (trinta) primeiros dias do semestre em curso. E devem ser considerados, ainda, os demais critérios de aproveitamento determinados no Titulo III, Capitulo IV, Seção I, do ROD, que trata do aproveitamento de componentes curriculares.

Já no que se refere à validação de conhecimentos, o IFCE validará conhecimentos adquiridos em estudos regulares ou em experiência profissional de estudantes do IFCE com situação de matrícula em matriculado, mediante avaliação teórica ou prática. O requerente poderá estar matriculado ou não no componente curricular para o qual pretende validar conhecimentos adquiridos. 
Não poderá ser solicitada validação de conhecimento para estágio curricular, trabalho de conclusão de curso e atividades complementares, assim como para estudantes que tenham sido reprovados no IFCE no componente curricular cuja validação de conhecimentos foi solicitada.

A solicitação de validação de conhecimentos deverá ser feita mediante requerimento protocolado e enviado à coordenadoria do curso, juntamente com o envio dos seguintes documentos: declaração, certificado ou diploma - para fins de validação em conhecimentos adquiridos em estudos regulares, cópia da Carteira de Trabalho (páginas já preenchidas) ou declaração do empregador ou de próprio punho, quando autônomo - para fins de validação de conhecimentos adquiridos em experiências profissionais anteriores e documentação complementar, caso seja solicitado pela comissão avaliadora.

O calendário do processo de validação de conhecimentos deverá ser instituído pelo próprio campus. Porém, a validação deverá ser solicitada nos primeiros 30 (trinta) dias do período letivo em curso e todo o processo de validação deverá ser concluído em até 50 (cinquenta) dias letivos do semestre vigente, a contar da data inicial de abertura do calendário do processo de validação de conhecimentos, definida pelo campus. 

A validação de conhecimentos de um componente curricular só poderá ser solicitada uma única vez e devem ser considerados, ainda, os demais critérios de aproveitamento determinados no Titulo III, Capitulo IV, Seção II, do ROD, que trata da validação de conhecimentos.
